\documentclass[11pt]{article}
\usepackage[margin=1in]{geometry}
\usepackage{booktabs,graphicx,amsmath,hyperref,xcolor,subcaption}
\usepackage[numbers]{natbib}
% find figures whether compiled from paper/ (../results) or the top-level (figures/)
\graphicspath{{../results/}{results/}{figures/}{./}}
\title{The Rendering Transformation Is Learnable; the Taxonomy Is Not:\\
A Dissociation in Synthetic--Real Sound-Effect Retrieval}
\author{SoundMatch-SR}
\date{Draft v2}
\begin{document}
\maketitle

\begin{abstract}
Perceptual studies report that listeners cannot reliably distinguish a well-made Foley or
text-to-audio rendering of an event from a real recording of the same event: the \emph{identity}
of the event survives the rendering. We ask whether modern audio embeddings represent that
identity invariantly to how a sound was produced, and find a sharp \textbf{dissociation}. We
introduce \textbf{SoundMatch-SR}, a benchmark for cross-domain (synthetic--real) sound-effect
retrieval, built on (i) the category-aligned DCASE-2023 Task-7 Foley corpus (5{,}550 real and
25{,}900 synthetic clips from 37 generator systems, 7 classes) and (ii) a new 34-category
UCS-labelled corpus (FSD50K real, CLAP-verified, with 10{,}420 instance-level Stable-Audio twins).
Off-the-shelf encoders carry a real, linearly-accessible gap (CLAP: $0.161$ mAP;
Proxy-A-distance $1.27$). A lightweight head on frozen embeddings closes it \emph{within} the
training taxonomy ($0.161\!\to\!0.018$), but a \emph{class-supervised} head does \emph{not}
transfer to unseen categories: held-out classes collapse below the frozen baseline, on both 7 and
34 categories. The fix is to train on the \emph{pairs}, not the labels: an instance-contrastive
objective (a clip and its generated twin) learns the rendering \emph{mapping} rather than category
identity, and generalizes---on categories never seen in training it retrieves the exact real twin
with R@1 $0.83\pm0.05$ (vs.\ $0.64$ frozen, $0.33$ class-supervised), across 5 folds and on three
encoders (CLAP, PANNs, AST). Category \emph{clustering} does not transfer for any objective. We
release the benchmark, the UCS corpus recipe, the trained heads, and a reusable transform.
\end{abstract}

\section{Introduction}
Foley research finds well-synthesized everyday sounds perceptually indistinguishable from
recordings of the same action; the event's identity survives the synthesis$\to$recording boundary
for a human listener. We make this measurable:
\emph{Is ``sound identity'' recoverable in embedding space, invariant to the rendering process---and
if so, does that invariance transfer to sounds the model was not trained on?}
We operationalize it as retrieval (detection rewards \emph{exploiting} the gap; retrieval rewards
\emph{removing} it) and answer with a dissociation that is, to our knowledge, new: the
synthetic$\to$real \emph{transformation} is a learnable, transferable object, while category
identity is not.

\paragraph{Contributions.}
\textbf{C1}: a two-part benchmark---a controlled 7-class corpus (DCASE-T7) and a diverse,
instance-paired 34-category UCS corpus (FSD50K + Stable Audio), with leakage-safe splits.
\textbf{C2}: a measured dissociation---an instance-contrastive objective generalizes the
synthetic$\to$real mapping to unseen categories; class-supervised invariance does not (and can
degrade unseen classes below baseline), robust over 5 folds and 3 encoders.
\textbf{C3}: a fidelity-spectrum analysis showing the instance head recovers the mapping across
twin-fidelity levels, collapsing only when the twin is independent of its source.
\textbf{C4}: released artifacts (benchmark, UCS recipe, CLAP-verification, trained heads, and a
reusable transform that adjusts any new clip's embedding).

\section{Related work}
General audio-embedding benchmarks (HEAR, X-ARES, MSEB) treat all audio as one domain and never
test invariance to \emph{how} a sound was produced. Deepfake-environmental-audio detection on
DCASE-T7 is the mirror objective (exploit the gap); we remove it and reproduce detection as a
controlled \emph{sensitive} head. We borrow Proxy-A-distance, DANN, Deep CORAL,
supervised-contrastive learning and IRM from domain adaptation, and connect to instance
discrimination: our generalizing objective is instance-level, and our negative result is that
class-level supervision does not transfer the invariance.

\section{The SoundMatch-SR benchmark}
\paragraph{Core corpus (DCASE-T7).} 7 classes; 5{,}550 real (DCASE source-disjoint dev/eval =
our train-val/test) and 25{,}900 synthetic clips from 37 generator systems. The controlled
setting for a clean gap measurement and a leave-one-class-out probe.
\paragraph{UCS corpus.} 34 UCS CatIDs mapping FSD50K's 200 AudioSet labels onto the
industry-standard SFX taxonomy. We CLAP-verify every clip (keep iff its label is top-5 by
audio--text cosine), retaining 10{,}420/13{,}579 ($77\%$). For each anchor we generate a Stable
Audio Open \emph{audio-init} twin conditioned on the real clip---10{,}420 instance pairs. Splits
follow FSD50K's train/val/eval, keyed on the Freesound uploader.
\paragraph{Tasks.} Category retrieval (relevant iff same category) with a same-domain control
$\to$ \emph{domain gap} = control $-$ cross; instance retrieval (relevant iff the exact
cross-domain twin) with R@1/MRR. Bootstrap CIs; on UCS, $k$-fold over held-out categories.

\section{Method}
A small MLP (CLAP-512$\to$256, $\ell_2$-normalised), frozen after training, applies to any new
clip. The objective decides what it learns: \textbf{invariant} (class-supervised contrastive,
cross-domain positives up-weighted, +DANN/IRM); \textbf{sensitive} (domain-supervised within
category---a fidelity axis); \textbf{instance} (positive = a clip and its own generated twin, with
pair-aware batching), which learns the synthetic$\to$real \emph{mapping}.

\section{Experiments}
\paragraph{Frozen encoders carry a real gap (Table~\ref{tab:frozen}).}
\paragraph{Within-taxonomy the gap closes; the two heads are mirror images.} The invariant head
closes the DCASE gap $0.161\!\to\!0.018$ while raising same-domain control ($0.935\!\to\!0.988$)
and making the space event-organised (silhouette by event $0.18\!\to\!0.75$, by domain
$0.04\!\to\!0.00$); the sensitive head flips it to organise by rendering (PAD $1.88$, silhouette by
domain $0.75$). Cross-domain supcon alone closes the gap; DANN/CORAL/IRM add $<0.002$.
\paragraph{The closed-world failure.} Holding one class out of training, the held-out class
collapses \emph{below} frozen (keyboard $0.69\!\to\!0.43$, gunshot $0.81\!\to\!0.31$, rain
$0.75\!\to\!0.30$). Adding categories does not fix it ($34$-cat: $0.40\!\to\!0.19$).
\paragraph{The dissociation (Table~\ref{tab:diss}, Fig.~\ref{fig:diss}).} Training on the pairs,
the instance objective generalizes the mapping to unseen categories (instance-R@1 $0.83$, every
fold), class-supcon hurts it ($0.33$), and no objective transfers category clustering. It holds on
PANNs and AST.
\paragraph{Not near-copies (Fig.~\ref{fig:fid}).} Across twin-fidelity levels the instance head
holds up (advantage $+0.43..0.52$) and collapses to chance only when the twin is independent of
its source.

\begin{table}[t]\centering
\caption{Frozen encoders: synthetic--real gap on DCASE-T7.}\label{tab:frozen}
\begin{tabular}{lccccc}\toprule
frozen & control & synth$\to$real mAP & gap & PAD & domain-probe\\\midrule
CLAP & 0.935 & 0.774\,[0.766,0.783] & $+0.161$ & 1.27 & 0.90\\
PANNs CNN14 & 0.798 & 0.674 & $+0.124$ & 1.39 & 0.90\\\bottomrule
\end{tabular}\end{table}

\begin{table}[t]\centering
\caption{5-fold leave-classes-out on UCS: held-out (unseen) categories, mean$\pm$std.}\label{tab:diss}
\begin{tabular}{lccc}\toprule
variant (unseen) & category-mAP & \textbf{instance-R@1} & instance-MRR\\\midrule
frozen & $0.615\pm0.052$ & $0.642\pm0.094$ & $0.731\pm0.081$\\
class-supcon & $0.459\pm0.072$ & $0.334\pm0.064$ & $0.439\pm0.061$\\
\textbf{instance} & $0.492\pm0.047$ & $\mathbf{0.829\pm0.045}$ & $0.888\pm0.034$\\\bottomrule
\end{tabular}\end{table}

\begin{figure}[t]\centering
\begin{subfigure}{0.49\textwidth}\includegraphics[width=\linewidth]{fig_dissociation.png}
\caption{}\label{fig:diss}\end{subfigure}\hfill
\begin{subfigure}{0.49\textwidth}\includegraphics[width=\linewidth]{fig_fidelity.png}
\caption{}\label{fig:fid}\end{subfigure}
\caption{(a) The instance/category dissociation on unseen categories, two encoders. (b) Fidelity
spectrum: the instance head holds up as twins diverge from their source.}
\end{figure}

\section{Limitations}
Category clustering does not transfer to unseen categories (train per taxonomy); twins are
Stable-Audio audio-init at one operating point per corpus; frozen-embedding heads cannot recover
information the backbone discarded; SSL encoders (BEATs/M2D/AudioMAE) are documented stubs.

\section{Release}
Code (MIT), DCASE + UCS manifests, CLAP-verification, generation recipe (not restricted audio),
trained heads + reusable transform, datasheet, and runbook.

\end{document}
